\documentclass[11pt]{article}
\usepackage{amsmath}
\usepackage[hidelinks]{hyperref}
\begin{document}
\title{Receipt, reservation and physical clearance across route replacement}
\author{}
\date{26 November 2026}
\maketitle

\paragraph{Audit.}
A current Open-RMF range, a previously delivered grant, and an actually
cleared crossing are separate inputs. On the pinned moderator, ordinary
route replacement installs a new reservation and resets the range.
The old readiness receipt does not authorize the new reservation.
After fresh readiness, the same numerical range can refer to a
different path. A physically disjoint new path can be used before the
old crossing clears, but reusing the old receipt for the new path
does not justify that choice. These are native grant observations in a
cross-simulator fixture; there is no result about a deployed robot's
actuation or an independent mission.

\paragraph{Interface.}
Pin the producer and reconstruct the old moderator from a durable
\texttt{SET}, \texttt{READY}, \texttt{REACHED} prefix before accepting
new work. The old participant is modeled as occupying crossing X.
The native authority source starts with the new participant waiting
at the near end; at tick 4 the old participant reports readiness
to depart. Under the geometry in this fixture, the range then grants
the second participant its first leg although the old participant
has not yet reached physical exit. The separately pinned warehouse
producer supplies concrete clearance timestamps 6, 18 and 19, with
clearance credited only after actual crate and carrier effects. The
bridge treats them as the old participant's departure times across
X and declares the clocks aligned. This does not derive Open-RMF
physical clearance from Open-RMF status or map a real warehouse's
layout to the Open-RMF geometry.

A normal \texttt{set} event at a chosen time replaces the second
participant's path. The replacement either keeps the same coordinates
or runs away from X along disjoint Y. Native \texttt{ready} on the
new reservation can reauthorize the new path. Redispatched readiness
for the old reservation is an ordinary delayed event in the controller
replay, not evidence of current authority. An independently booked
route reaches its checkpoint at tick 13 in every replay. The second
crossing must finish by tick 20, the carried clearance by 19, and
the completion remains a modeled one-tick effect.
The driver retains actual ordered tuples of reservation identity,
granted end, geometry and incumbent status; it does not replace the
native decision with a Boolean from the checker.

\paragraph{Two authority laws.}
An option carries its observed grant identity and exact path,
as well as the chosen launch time. A receipt at tick 5 and, when
available, a new observation at tick 11 can qualify that option.
The interpreter checks actual physical clearance, precise active
grant, and the carried and independent effects. The second
observation cannot establish what a later publisher will do.
No choice sees the private warehouse completion time. Scenario
classes include both early and late cases, conditioning only on
the declared readiness schedule and the observation rule.

Two downstream laws are deliberately kept distinct. Under
\emph{issue-only authority}, issuing an approved command authorizes
its modeled effect to finish; subsequent edits do not retroactively
rename the effect. Under \emph{through-effect authority}, the
same authorization must extend through modeled completion too.
The pinned moderator identifies the active grant, not which of these
plant contracts holds. With a single grant cutoff $r$ and ties
ordered as update before use, an issue at $t$ with duration $d$
separates the laws exactly when $t<r\leq t+d$. For arbitrary
replacement, eligibility is also identity- and geometry-indexed.

\paragraph{Witnesses and limits.}
At the same disclosed view at tick 11, a permitted old operation at
tick 12 can be admissible when the change is at 13, and inadmissible
when it is at 12. Their exact common choice set is empty in the
selected two histories. If the change occurs after issue but before
completion, the case separates the two authority laws; an operation
under the issue-only law does not consume a fresh permit for an
unexecuted action. When the edit comes earlier and current Y is
requalified, it can finish while X remains physically occupied.
Exactly the same coordinates under a new reservation still require
authority for the new version: they preserve geometry, not identity.

A different interface can requalify inline at the acting source.
Five selected replays keep the same moderator process live through
the modeled command and issue a native reached event only if the
current grant still matches at completion. One run selects X and one
selects newly granted disjoint Y; missing current authority or live X
blocks issue. A replacement within the modeled execution interval
prevents the old acknowledgement under the through-effect law. All
five preserve the unrelated reached event. This is a serialized
local-source control, not a claim that a remote snapshot is a lease,
that a physical motion can be canceled instantly, or that a native
status alone verifies the modeled progress.

The tests cover 240 native moderator schedules, 5,040 logged
native cuts, ten warehouse runs, and 11,520 whole-trajectory
option comparisons against direct checks of both native authority
and physical clearance. All direct comparisons agree. Grouping
without revealing the physical future gives 32 local view classes
across the enumerated families; four admit a common option. These
are development counts, not probabilities, and do not measure
planner superiority. Lean checks the time boundary and the
identity/common-view obstructions, not the native code. The setup
does not solve stable handoff, reauthorization under a new site
controller, or safety after a change that was not included in the
admitted support.

An appropriate comparison cannot assume a global pause. Existing
\href{https://whoenig.github.io/publications/2019_RA-L_Hoenig.pdf}{action
dependency graph execution} already accounts for relevant predecessor
completion while continuing disjoint work. The strongest stock
baseline here begins from the durable reconstruction and is given
the same scoped observations. The remaining test is a frozen selector
facing independently supplied authority and physical constraints, not
interpreting the absence of a new message as evidence that a grant
stayed in force.
\end{document}
